\documentclass[10pt,a4paper]{amsart}
\usepackage{amsmath,amssymb,amsthm,graphicx}
\usepackage{hyperref}
\newcommand{\C}{\mathbb{C}}
\renewcommand{\H}{\mathbb{H}}
\newcommand{\Z}{\mathbb{Z}}
\newcommand{\T}{\mathbb{T}}
\newcommand{\Q}{\mathbb{Q}}
\newcommand{\N}{\mathbb{N}}
\newcommand{\R}{\mathbb{R}}
\newcommand{\slz}{{\text {\rm SL}}_2(\mathbb{Z})}
\newcommand{\re}{\textnormal{Re}}
\newcommand{\im}{\textnormal{Im}}
\newcommand{\Tr}{\operatorname{Tr}}

\newcommand{\vt}[1]{\left\lvert #1 \right\rvert}
\newcommand{\rf}{\mathfrak{r}}
\newcommand{\Ec}{\mathcal{E}}
\newcommand{\Cc}{\mathcal{C}}
\newcommand{\Qc}{\mathcal{Q}}
\newcommand{\Fc}{\mathcal{F}}

\numberwithin{equation}{section}

\newtheorem{Theorem}{Theorem}[section]
\newtheorem*{Theorem*}{Theorem}
\newtheorem{Corollary}[Theorem]{Corollary}
\newtheorem{Lemma}[Theorem]{Lemma}
\newtheorem{Proposition}[Theorem]{Proposition}
{ \theoremstyle{definition}
	\newtheorem{Definition}[Theorem]{Definition}
	\newtheorem{Note}[Theorem]{Note}
	\newtheorem{Example}[Theorem]{Example}
	\newtheorem{Remark}[Theorem]{Remark} }
\usepackage{cleveref}
\crefname{originalreference}{original reference}{original references}
\Crefname{originalreference}{Original reference}{Original references}
\makeatletter
\newlabel{Prop: Craig}{{Prop--Craig}{source}{}{sourcepointer-1}{}}
\newlabel{Prop: Craig@cref}{{[originalreference][0][]Prop--Craig}{[1][0][]source}{}{}{}}
\makeatother
\title{Equidistribution of the three partition statistics in residue classes: the three asymptotic formulas of Theorem 1.1}
\author{Males, Joshua}
\date{doi:10.3842/SIGMA.2024.001; CC BY 4.0}
\begin{document}
\maketitle
\setcounter{section}{1}
\setcounter{Theorem}{0}
\noindent\textit{Editorial note.} Author's own Theorem environment, counted statement at source lines 107-125 and its own proof environment at source lines 310-341, both read verbatim from sources/sigma-2024-001-a-note-on-the-equidistribution-of-3-colour-p/pinned-original-source.tex. No line of either span was dropped or reworded; in particular the proof is delivered with its optional argument intact.\par\medskip
\noindent\textit{Editorial note.} Detached proof, recorded on the card face: the statement is in the introduction and the proof stands in its own section, headed Proof of Theorem 1.1, more than a hundred source lines later. The proof is counted here under the position-independent ruling for detached proofs, with the content of the span, not its position, deciding the count. The census does not pair this statement with any proof block, so the pairing on this card is read from the source text.\par\medskip
\noindent\textit{Editorial note.} Work delegated inside the counted span, recorded verbatim so that no reader takes it for a reproduced argument: the minor-arc estimate is taken from the published arguments of the reference Zhou with the parameter choices spelled out in the text, and the proof states that its minor-arc treatment is very similar to that of the reference BCMO; the discussion of the remaining cases ends with the sentence that the author leaves the details of those further cases to the interested reader.\par\medskip
\noindent\textit{Editorial note.} The article is distributed under CC BY 4.0 with the fixed version URL https://sigma-journal.com/2024/001/ and the licence observation is filed next to the source as license-observation.json. The counter seed is taken from the number this Theorem carries in the published article, 1.1 (census version 3.7 field printed), which is the article's own first numbered theorem. The unit was scanned against the published-item list by DOI before build: no published entry carries this DOI. All card fields are ASCII only.\par\medskip
\noindent\textit{Allowed standard prerequisites (editorial).}\begin{itemize}
\item Generating functions of the three statistics as products of F\_\{a,c\}. Same article, source lines 286-290, opening of the section that proves Theorem 1.1, not delivered on this card. Supplies the three identities that write J\_E, J\_T and J\_G as products of F\_\{a,c\}(zeta;q) and one q-Pochhammer product, which is what the proof takes its asymptotics of.
\item Asymptotics of the reciprocal q-Pochhammer product. Same article, the displayed well-known bound at source lines 292-295, not delivered on this card. Supplies the expansion of (e\textasciicircum{}\{-z\};e\textasciicircum{}\{-z\})\_infty\textasciicircum{}\{-1\} used to pass from F\_\{a,c\} to each J.
\item Asymptotics of the J at roots of unity. Same article, the Corollary at source lines 297-306, not delivered on this card. Supplies the three asymptotic formulas for J\_E, J\_T and J\_G at a primitive root of unity, obtained there by combining the two items above with the asymptotics of F\_\{a,c\}.
\item Asymptotics of F\_\{a,c\} at zeta equal to 1. Same article, the Proposition quoted from the reference Craig at source lines 276-282, registered on this card as a non-load-bearing navigation pointer. Supplies the asymptotic formula for F\_\{a,c\}(1;e\textasciicircum{}\{-z\}) that the major-arc computation invokes when the multisection's j = 0 term is evaluated.
\item Domination of the pole at q = 1. Published result cited in the proof as the reference Chern, Theorem 1.1 of that paper. Supplies the statement that for each of the infinite products the pole at q = 1 dominates the poles at the other roots of unity, which the proof invokes in its discussion of the remaining cases.
\end{itemize}\medskip
\ifdefined\hypertarget\hypertarget{sourcepointer-1}{}\fi
\noindent\textit{Source reference.} Original reference Prop--Craig: sources/sigma-2024-001-a-note-on-the-equidistribution-of-3-colour-p/pinned-original-source.tex, source lines 276--276. Navigation pointer invoked at proof line 321; the counted span contains the pointer only and the cited result is not reproduced inside it, so the pointer is registered as a non-load-bearing navigation stub in the established form for Stacks material. No bibliography citation is involved.\par\medskip
\begin{Theorem}\label{thm: main}
	The following equidistribution results are true.
	\begin{enumerate}
		\item[$(1)$] Let $s \in \N$ such that $3 \nmid s$. Then
		\begin{align*}
			J_E(r,s;n) = \frac{1}{4 \sqrt{3} s n}{\rm e}^{\pi \sqrt{\frac{2n}{3}}} \big(1+O\big(n^{-\frac{1}{2}}\big)\big).
		\end{align*}
		
		\item[$(2)$] Let $s \in \N$. Then
		\begin{align*}
			J_T(r,s;n) = \frac{1}{2^{\frac{7}{2}} s n^{\frac{3}{2}}} {\rm e}^{\pi\sqrt{2n}} \big(1+O\big(n^{-\frac{1}{2}}\big)\big).
		\end{align*}
		
		\item[$(3)$] Let $s\in \N$ such that $2 \nmid s$. Then
		\begin{align*}
			J_G(r,s;n) = \frac{{\rm e}^{\pi \sqrt{n}}}{8sn} \big(1+O\big(n^{-\frac{1}{2}}\big)\big).
		\end{align*}
	\end{enumerate}
\end{Theorem}

\begin{proof}[Proof of Theorem \ref{thm: main}]
It is a simple exercise using orthogonality of roots of unity to show that
\begin{align}\label{eqn: multisection}
H_*(r,s;q) := \sum_{n \geq 0} J_*(r,s;n) q^n = \frac{1}{s} \sum_{j=0}^{s-1} \zeta_b^{-rj} J_*\big(\zeta_s^j;q\big).
\end{align}

We now use Wright's circle method to show that the $j=0$ term dominates the asmyptotics for each choice $* \in \{ E,T,G \}$, which in turn implies that the main term asymptotic is independent of the residue class and hence the statistic is asymptotically equidistributed.

One may be concerned that poles other than $q=1$ may contribute to the main asymptotic term. However, by Theorem~1.1 of the beautiful paper of Chern~\cite{Chern} we know that for each of the infinite products, the pole at $q=1$ dominates, since each $F_{a,c}^{-1}$ is exponentially smaller at poles other than $q=1$. This is only true because of the assumptions on $s$ made in each case. For example, in case (3) if one had terms where $\zeta = -1$ then the pole at $q =-1$ would also have a term with the same growth as the main term asymptotic, and would cancel the main term already obtained thanks to the sum over $j$ in \eqref{eqn: multisection}. We leave the details of these further cases to the interested reader.


Using that $\lvert \zeta \Phi(\zeta,2,1) \rvert < \frac{\pi^2}{6}$ for $\zeta$ is a primitive root of unity not equal to~$1$, it is simple to see that on the major arc the $j=0$ term in \eqref{eqn: multisection} exponentially dominates all other terms, and thus using Proposition~\ref{Prop: Craig}, we have
\begin{align*}
		&H_E(r,s;q) = \frac{\Gamma\left(\frac{1}{3}\right) \Gamma\left(\frac{2}{3}\right)}{s (2\pi)^{\frac{3}{2}}} \sqrt{3z} \exp\left( \frac{\pi^2}{6z} \right) \left(1+O(|z|)\right) , \\
	&H_T(r,s;q) = \frac{1}{s (2\pi)^{\frac{3}{2}}} z^{\frac{3}{2}} \exp\left(\frac{\pi^2}{2z}\right) \left(1+O(|z|)\right), \\
	&H_G(r,s;q) = \frac{1}{2s} \sqrt{\frac{z}{\pi}} \exp\left( \frac{\pi^2}{4z} \right)\left(1+O(|z|)\right).
\end{align*}
Note that by Euler's reflection formula for the $\Gamma$-function, we obtain $\Gamma\left(\frac{1}{3}\right) \Gamma\left(\frac{2}{3}\right) = \frac{2 \pi}{\sqrt{3}}$.


 It remains to estimate each of $H_*(r,s;q)$ on the minor arcs, where $Mx \le |y|\le \pi$. To do so, we estimate $F_{a,c}(\zeta;q)$ on the minor arcs using arguments very similar to those given in \cite{BCMO} along with the classical fact that for some $\mathcal{C}>0$, we have
 \begin{align*}
 \big|({\rm e}^{-z};{\rm e}^{-z})_\infty^{-1}\big| \le x^{\frac12} {\rm e}^{\frac{\pi}{6x}-\frac{\mathcal{C}}{x}}.
 \end{align*}
This guarantees that the final $q$-Pochhammer in each case is exponentially smaller on the minor arc, and so we need only consider the $F_{a,c}$ functions.

We aim to show that (for $Mx \le |y|\le \pi$)
\begin{align*}
	\re\big(\log F_{a,c}^{-1}\big(\zeta_b^j ; q \big) \big) - \log\big(F_{a,c}^{-1}(1;|q|)\big) \leq - \frac{\mathcal{C}}{x}
\end{align*}
for some positive constant $\mathcal{C}$. This follows from the arguments of \cite[p.~1962]{Zhou} (see, in particular, line~7) with the parameters $f(n) = cn+ (a-c)$, $k=b$, $\delta = 1$, and $\ell =0$. Therefore, the contribution from the major arc exponentially dominates, and we obtain asymptotic equidistribution as claimed.
\end{proof}

\begin{thebibliography}{99}
\footnotesize\itemsep=0pt

\bibitem{BCMO}
Bringmann K., Craig W., Males J., Ono K., Distributions on partitions arising
 from {H}ilbert schemes and hook lengths, \href{https://doi.org/10.1017/fms.2022.45}{\textit{Forum Math. Sigma}}
 \textbf{10} (2022), e49, 30~pages, \href{https://arxiv.org/abs/2109.10394}{arXiv:2109.10394}.

\bibitem{Campbell}
Campbell R., Les int\'egrales eul\'eriennes et leurs applications. \'{E}tude
 approfondie de la fonction gamma, \textit{Collection Universitaire de
 Math\'ematiques}, Vol.~20, Dunod, Paris, 1966.

\bibitem{Chern}
Chern S., Nonmodular infinite products and a conjecture of {S}eo and {Y}ee,
 \href{https://doi.org/10.1016/j.aim.2023.108932}{\textit{Adv. Math.}} \textbf{417} (2023), 108932, 40~pages,
 \href{https://arxiv.org/abs/1912.10341}{arXiv:1912.10341}.

\bibitem{Ciolan}
Ciolan A., Equidistribution and inequalities for partitions into powers,
 \textit{Bull. Math. Soc. Sci. Math. Roumanie~(N.S.)} \textbf{66} (2023),
 409--431, \href{https://arxiv.org/abs/2002.05682}{arXiv:2002.05682}.

\bibitem{C}
Craig W., Seaweed algebras and the index statistic for partitions,
 \href{https://doi.org/10.1016/j.jmaa.2023.127544}{\textit{J.~Math. Anal. Appl.}} \textbf{528} (2023), 127544, 23~pages,
 \href{https://arxiv.org/abs/2112.09269}{arXiv:2112.09269}.

\bibitem{HardyRamanujan}
Hardy G.H., Ramanujan S., Asymptotic formulae in combinatory analysis,
 \href{https://doi.org/10.1112/plms/s2-17.1.75}{\textit{Proc. Lond. Math. Soc.}} \textbf{17} (1918), 75--115.

\bibitem{LZ}
Liu Z.-G., Zhou N.H., Uniform asymptotic formulas for the {F}ourier coefficients
 of the inverse of theta functions, \href{https://doi.org/10.1007/s11139-021-00409-8}{\textit{Ramanujan~J.}} \textbf{57} (2022),
 1085--1123, \href{https://arxiv.org/abs/1903.00835}{arXiv:1903.00835}.

\bibitem{NgoRhoades}
Ngo H.T., Rhoades R.C., Integer partitions, probabilities and quantum modular
 forms, \href{https://doi.org/10.1186/s40687-017-0102-4}{\textit{Res. Math. Sci.}} \textbf{4} (2017), 17, 36~pages.

\bibitem{SZ}
Schlosser M.J., Zhou N.H., Expansions of averaged truncations of basic
 hypergeometric series, \href{https://arxiv.org/abs/2307.10821}{arXiv:2307.10821}.

\bibitem{Wright1}
Wright E.M., Stacks, \href{https://doi.org/10.1093/qmath/19.1.313}{\textit{Quart.~J. Math. Oxford Ser.~(2)}} \textbf{19}
 (1968), 313--320.

\bibitem{Wright2}
Wright E.M., Stacks.~{II}, \href{https://doi.org/10.1093/qmath/22.1.107}{\textit{Quart.~J. Math. Oxford Ser. (2)}} \textbf{22}
 (1971), 107--116.

\bibitem{Zhou}
Zhou N.H., Note on partitions into polynomials with number of parts in an
 arithmetic progression, \href{https://doi.org/10.1142/S1793042121500718}{\textit{Int.~J. Number Theory}} \textbf{17} (2021),
 1951--1963, \href{https://arxiv.org/abs/1909.13549}{arXiv:1909.13549}.

\end{thebibliography}
\end{document}
